% Chapter 5 -- The File Commands (AN-08, AN-09).
% Sources (Law 1): os/i123/app/face_a.c paint_command (the FILES prompts
% "Name of file to retrieve: " / "Enter name of file to save: ", the
% default A:\*.wk1 and the current file), file_name (the .wk1 default
% extension; "Enter a file name", "Invalid file name"), retrieve (the
% transaction: a failed Retrieve leaves the worksheet; "File not found",
% "Path not found", "File I/O error", "File too large", "File read
% incomplete", "Not enough worksheet memory", "Not a valid WK1 file"),
% save / write_sheet (the Cancel / Replace / Backup menu, its
% descriptions, Backup refusing; the status line "Saved ..." /
% "Retrieved ..."; "Worksheet too large to save"), os/i123/fs/wk1.c (the
% .WK1 codec); spec/i123/era_shortfalls.md I123-P3-02 (the other File
% leaves refuse).  Menu tree generated from the shipped resource.
% GATES: test-i123-wk1-roundtrip test-flair-i123-files test-flair-i123-file-ui test-i123-menu-tree test-pubs-i123 test-pubs-screens
\Gates{test-i123-wk1-roundtrip, test-flair-i123-files, test-flair-i123-file-ui,
  test-i123-menu-tree, test-pubs-i123, test-pubs-screens}
\chapter[The File Commands]{\textit{The File Commands}}
\input{i123-tree-File}

The File commands store worksheets on disk and bring them back into
memory.  A worksheet you have not saved exists only in memory, and is lost
when you end Initech 123 or erase the worksheet.

\begin{LotTable}{p{30mm}p{98mm}}
\LotHead{\textbf{Command} & \textbf{Task}}
/File Retrieve & Replaces the worksheet in memory with a worksheet file\\
/File Save & Saves the worksheet in memory in a worksheet file\\
/File Combine & Opens the menu for incorporating another file \draft{I123-P3-02: its commands refuse}\\
/File Xtract & Opens the menu for saving part of a worksheet \draft{I123-P3-02: its commands refuse}\\
/File Erase & Opens the menu for erasing files \draft{I123-P3-02: its commands refuse}\\
/File List & Opens the menu for listing files \draft{I123-P3-02: its commands refuse}\\
/File Import & Opens the menu for reading text files \draft{I123-P3-02: its commands refuse}\\
/File Directory & Displays or changes the current directory \draft{I123-P3-02: refuses}\\
/File Admin & Opens the menu of file administration \draft{I123-P3-02: its commands refuse}\LotEnd
\end{LotTable}

\section{Common Uses of the File Commands}
\begin{itemize}
\item Saving your work (\slashcmd{/File Save})
\item Saving a second copy of a worksheet under another name
  (\slashcmd{/File Save})
\item Continuing work on a worksheet you saved earlier
  (\slashcmd{/File Retrieve})
\end{itemize}

\section{Related Reading}
\begin{itemize}
\item ``Worksheet Files'' in Chapter~1
\item Chapter~9, ``The Quit Command'', on saving your work before you end
  a session
\item Lessons~6, 7, 11 and 17 in the \doctitle{Tutorial}
\end{itemize}

\section{Terms You Need to Know}
\begin{itemize}
\item A \term{worksheet file} is a file on disk that holds a worksheet: its
  entries, the formats and protection settings of its cells, its column
  widths, its range names and its global settings.
\item A \term{file name} has up to eight characters, optionally followed by
  a period and an extension of up to three characters.  Worksheet files
  have the extension \file{.WK1}.
\item The \term{current file} is the file you last retrieved or saved.
  /File Save offers its name.
\end{itemize}

\section{Before You Use the File Commands}
Initech 123 reads and writes files on the InitechOS data disk, drive
\file{A:}.  A file name you type without a drive refers to drive \file{A:}.
You can type file names in capital or small letters.

\begin{Caution}
/File Retrieve replaces the worksheet in memory.  Save the worksheet you
are working on first if you want to keep it.
\end{Caution}

\MenuEntry{/File}
/File displays the File menu.  This version carries out /File Retrieve and
/File Save; the other File commands appear in the menu and their own menus
can be displayed, but they are not available in this version.

\MenuEntry{/File Retrieve}
/File Retrieve reads a worksheet file into memory, replacing the worksheet
that was there.
\FlairShot[fig:retrieve]{retrieve}{The Retrieve prompt}

\BHead{Terms You Need to Know}
\begin{itemize}
\item The \term{current file} becomes the file you retrieve.
\end{itemize}
Initech 123 reads the file completely and checks it before it changes
anything.  Only when the whole file has been read as a valid \file{.WK1}
worksheet does the retrieved worksheet replace the one in memory, with its
entries, formats, column widths, range names and global settings, and the
cell pointer where the file left it.  If anything goes wrong, the worksheet
in memory is left exactly as it was.

\BHead{Procedure}
\begin{Steps123}
\item Select /File Retrieve.  The mode indicator changes to \ind{FILES},
  and the control panel shows \texttt{Name of file to retrieve:} and
  \texttt{A:\textbackslash *.wk1}.
\item Type the name of the file and press \kname{ENTER}.

  You may leave out the extension: Initech 123 adds \file{.WK1}.  The
  status line shows \texttt{Retrieved} and the name of the file, and
  Initech 123 returns to \ind{READY} mode.
\end{Steps123}

\begin{Note}
If Initech 123 cannot retrieve the file, the mode indicator changes to
\ind{ERROR} and the control panel shows the reason; press \kname{ESC} to
return to the worksheet.  The messages are:
\begin{LotTable}{p{44mm}p{84mm}}
\LotHead{\textbf{Message} & \textbf{Meaning}}
\Msg{File not found} & There is no file of that name\\
\Msg{Path not found} & A directory named in the file name does not
  exist\\
\Msg{Not a valid WK1 file} & The file is not a worksheet file, or it is
  damaged\\
\Msg{File too large} & The file is larger than Initech 123 can read
  (64K)\\
\Msg{Not enough worksheet memory} & The worksheet in the file has more
  entries than Initech 123 can hold (Appendix~A)\\
\Msg{File read incomplete} & The file could not be read to its end\\
\Msg{File I/O error} & The disk could not be read\\
\Msg{Enter a file name} & You pressed \kname{ENTER} without typing a
  name\\
\Msg{Invalid file name} & The name contains a space, \texttt{*} or
  \texttt{?}\\
\Msg{File name too long} & You typed more than 48 characters\LotEnd
\end{LotTable}
\end{Note}

\MenuEntry{/File Save}
/File Save saves the worksheet in memory, complete, in a worksheet file.
\FlairShot[fig:saveprompt]{saveprompt}{The Save prompt}

\BHead{Procedure}
\begin{Steps123}
\item Select /File Save.  The control panel shows
  \texttt{Enter name of file to save:} followed by the name of the current
  file, or by \texttt{A:\textbackslash *.wk1} if the worksheet has not yet
  been saved or retrieved.
\item Press \kname{ENTER} to save under the name shown, or type a new name
  and press \kname{ENTER}.

  If there is no file of that name, Initech 123 saves the worksheet, shows
  \texttt{Saved} and the file name on the status line, and returns to
  \ind{READY} mode.
\item If a file of that name already exists, Initech 123 displays a menu of
  three choices (Figure~\ref{fig:replace}).  Select one:
\begin{Options}
\item[\normalfont\bfseries Cancel\hfill] Leaves the file on disk as it is;
  Initech 123 returns to \ind{READY} mode without saving.
\item[\normalfont\bfseries Replace\hfill] Replaces the file on disk with
  the worksheet in memory.
\item[\normalfont\bfseries Backup\hfill] \draft{not available in this
  version: displays Command not implemented} Renames the file on disk with
  the extension \file{.BAK} and then saves the worksheet.
\end{Options}
\end{Steps123}

\FlairShot[fig:replace]{replace}{Saving under the name of an existing file}

\begin{Note}
Initech 123 checks that the worksheet can be written before it creates or
replaces the file.  If it cannot (\Msg{Worksheet too large to save}), the
file on disk is not touched.  If the disk fails while a new file is being
written, Initech 123 removes the incomplete file and displays
\Msg{File I/O error}.
\end{Note}

\LabelHead{Tip}
\begin{itemize}
\item Save often.  A saved worksheet is safe from a mistaken
  \slashcmd{/Worksheet Erase Yes} and from a power failure.
\item To keep the previous version of a worksheet, save the changed
  worksheet under a new name.
\end{itemize}

\MenuEntry{/File Combine}
/File Combine displays the menu for incorporating all or part of a
worksheet file into the worksheet in memory.
\begin{Options}
\OptionMenu{Copy} Displays the menu for copying data from the file over the
  worksheet.
\OptionMenu{Add} Displays the menu for adding the file's values to the
  worksheet's.
\OptionMenu{Subtract} Displays the menu for subtracting the file's values from
  the worksheet's.
\end{Options}
\SubEntry{/File Combine Copy}
\begin{Options}
\Option[I123-P3-02: refuses]{Entire-File} Copies a whole file into the
  worksheet at the cell pointer.
\Option[I123-P3-02: refuses]{Named/Specified-Range} Copies a range of a
  file into the worksheet at the cell pointer.
\end{Options}
\SubEntry{/File Combine Add}
\begin{Options}
\Option[I123-P3-02: refuses]{Entire-File} Adds the values of a whole file
  to the worksheet.
\Option[I123-P3-02: refuses]{Named/Specified-Range} Adds the values of a
  range of a file to the worksheet.
\end{Options}
\SubEntry{/File Combine Subtract}
\begin{Options}
\Option[I123-P3-02: refuses]{Entire-File} Subtracts the values of a whole
  file from the worksheet.
\Option[I123-P3-02: refuses]{Named/Specified-Range} Subtracts the values
  of a range of a file from the worksheet.
\end{Options}

\MenuEntry{/File Xtract}
/File Xtract displays the menu for saving part of the worksheet in a file of
its own.
\begin{Options}
\Option[I123-P3-02: refuses]{Formulas} Saves a range with its formulas.
\Option[I123-P3-02: refuses]{Values} Saves a range with the current values
  of its formulas.
\end{Options}

\MenuEntry{/File Erase}
/File Erase displays the menu for erasing a file from disk.
\begin{Options}
\Option[I123-P3-02: refuses]{Worksheet} Erases a worksheet file.
\Option[I123-P3-02: refuses]{Print} Erases a print file.
\Option[I123-P3-02: refuses]{Graph} Erases a graph file.
\Option[I123-P3-02: refuses]{Other} Erases a file of any kind.
\end{Options}

\MenuEntry{/File List}
/File List displays the menu for listing the files on disk.
\begin{Options}
\Option[I123-P3-02: refuses]{Worksheet} Lists the worksheet files.
\Option[I123-P3-02: refuses]{Print} Lists the print files.
\Option[I123-P3-02: refuses]{Graph} Lists the graph files.
\Option[I123-P3-02: refuses]{Other} Lists every file.
\Option[I123-P3-02: refuses]{Linked} Lists the files the worksheet's
  formulas refer to.
\end{Options}

\MenuEntry{/File Import}
/File Import displays the menu for reading a text file into the worksheet.
\begin{Options}
\Option[I123-P3-02: refuses]{Text} Reads each line of the file into a cell
  as a label.
\Option[I123-P3-02: refuses]{Numbers} Reads the numbers and quoted text of
  each line into separate cells.
\end{Options}

\MenuEntry[I123-P3-02: refuses]{/File Directory}
Displays the current directory and lets you change it.

\MenuEntry{/File Admin}
/File Admin displays the menu of file administration commands.
\begin{Options}
\OptionMenu{Reservation} Displays the Reservation menu.
\Option[I123-P3-02: refuses]{Table} Enters a table of file information in
  the worksheet.
\Option[I123-P3-02: refuses]{Link-Refresh} Updates the formulas that refer
  to other files.
\end{Options}
\SubEntry{/File Admin Reservation}
\begin{Options}
\Option[I123-P3-02: refuses]{Get} Obtains the reservation that permits
  saving a shared file.
\Option[I123-P3-02: refuses]{Release} Gives up the reservation.
\end{Options}
